\cotefile{157-5}
\scanpage{261}\mspage{177}
% archive sheet number: 260
\[
\left\{\begin{aligned}
M_{\pi}&\subset\mathcal F^{\unsure{\pi|P}}_{\deleted{\unsure{\pi}}}
  &&(\text{si }\pi=\pi^{\pi_0}),\\
\text{donc }\ M_{\pi}&=\mathcal F_{\pi}
  &&\text{si \inserted{de plus} la \unsure{fcg|fag} est \uline{accessible}.}
\end{aligned}\right.
\tag*{(4.7.9.1)}
\]
\deleted{\illegible{1}}

\uline{Dém} \unsure{Appliquant|Appliquons} la condition c) des fcg à
un $F\in\operatorname{Ob}M_\pi$, on trouve que $\deleted{\unsure{F}}F$ est facteur
direct d'un $\deleted{\unsure{E}}F'\in\operatorname{Ob}\mathcal F_\pi$. Comme $\mathcal F^\pi_{\deleted{\unsure{\pi}}}$ est
stable par $\varinjlim$ \unsure{grdes|gdes} devant $\pi_0$ et de cardinal
$\leq\pi$, et que \uline{de plus} $\pi>\pi_0$, on conclut que

\uline{Lemme 4.7.9.1} : Si $\pi>\pi_0$,$^{*}$ $\mathcal F^\pi_{\deleted{\unsure{\pi}}}$ est stable par
facteurs directs,

% slanted
\margin{Mais $\mathcal F^{\pi_0}$ n'est pas forcément stable par facteurs directs, et comme on veut \unsure{retrouver} $C$ prenant $M=\operatorname{Ind}_{\pi_0}(C)$, $C$ une petite catégorie, on \unsure{ne voudra|voudrait} \unsure{pas|prendre} \unsure{ds} la filtration \unsure{cardinale|considérée}, \unsure{prendre} les $\mathrm{Filt}'^{\pi}$ --- ou, \unsure{notons} que $\mathrm{Filt}'^{\pi_0}=\mathrm{Filt}^{\pi_0}$ \unsure{S.} $\unsure{I\simeq C}$\deleted{\unsure{est tout petit}}, ou \unsure{quiconque|toute} cat. \deleted{\unsure{petite}} \unsure{grde} devant $\pi_0$ et de cardinal $\leq\pi_0$ est \unsure{cofinale} à $e$.}
il résulte que $F\in\mathcal F^\pi$, d'où (4.7.9.1).

\uline{Corollaire} \deleted{\unsure{4.}}4.7.10. Supposons la fcg accessible.
Pour tout cardinal $\pi\geq\pi_0$, on a
\[
\mathcal F^\pi_{\deleted{\unsure{\pi}}}\subset M_\pi\subset\mathcal F^{\pi^{\pi_0}}
\quad(=M^{\unsure{P|\pi}}_{\pi^{\pi_0}})
\tag*{(4.7.10.1)}
\]

En effet, $\mathcal F_\pi\subset M_\pi$ par définition de fcg
accessible, et $M_\pi\subset\deleted{\mathcal F^{\pi^{\pi_0}}}M_{\pi^{\pi_0}}$ (puisque
$\pi\leq\pi^{\pi_0}$), or $M_{\pi^{\pi_0}}=\mathcal F^{(\pi^{\pi_0})}$ en vertu de
(4.7.9.1) appliqué à $\pi'=\pi^{\pi_0}$ \unsure{ald} $\pi$ (vu que
$\pi'^{\pi_0}=\pi'$), d'où (4.7.10.1). On notera
que cette inclusion généralise (4.7.9.1),
elle donne $\mathcal F^\pi_{\deleted{\unsure{\pi}}}=M_\pi$ si $\pi=\pi^{\pi_0}$.

4.7.11. \deleted{\unsure{Il}} Il \inserted{me} semble que les seules \deleted{\unsure{condi}} filtrations,
cardinales \inserted{\unsure{gén.}} utiles sont celles qui sont \uline{accessibles}.
\deleted{\unsure{Et}} Dans la suite je vais me borner à ce
cas. Comme \inserted{alors} $\mathcal F_\pi\subset M_\pi$, on en conclut
\scanpage{262}\mspage{\unsure{178}}
% archive sheet number: 261
que le foncteur canonique
\[
\operatorname{Ind}_{\pi}(\mathcal F^\pi_{\deleted{\unsure{\pi}}})\hookrightarrow M
\tag*{(4.7.11.1)}
\]
prolongeant l'inclusion $\mathcal F^\pi\subset M$, de façon
à commuter aux $\varinjlim$ \unsure{grdes} devant $\pi$, est
\uline{pleinement} \uline{fidèle}. De plus, lorsque
$\pi^{\pi_0}=\pi$, la condition c) des \unsure{fcg|fag} nous dit
\inserted{entre autres} que ce foncteur est ess. surjectif, donc
que c'est une \uline{équivalence de catégories} :
\[
\operatorname{Ind}_{\pi}(\mathcal F^\pi)\xrightarrow{\approx}M_{\unsure{\sim|a}}\quad
(\text{si }\pi=\pi^{\pi_0}).
\tag*{(4.7.11.2)}
\]

\uline{Proposition} 4.7.12. Soit $(\mathcal F^\pi)_{\pi\geq\pi_0}$ une fcg
\inserted{accessible (\unsure{fcga})} sur la $\unsure{U|M}$-cat. $M$, et soit $\pi$ un cardinal $\geq\pi_0$.
Les conditions suivantes sont équivalentes

a) Le foncteur pl. fidèle (4.7.11.1) est
ess. surjectif, i.e. est une équivalence de
catégories.

b) Pour \unsure{tout} $F\in\operatorname{Ob}M$, $\mathcal F^\pi_{/F}$ est
une catégorie \inserted{filtrante et} grande devant $\pi$.

\deleted{\unsure{On}} Quand ces conditions sont satisfaites, \inserted{et si $\pi>\pi_0$} on a
\deleted{Résulte aussitôt de 4.4.4.}
\[
\mathcal F^\pi=M_\pi\mathpunct{.}
\tag*{(4.7.12.1)}
\]

\uline{Dém.} L'équivalence de a) et b) résulte
aussitôt de 4.4.\unsure{2|8}. et de la définition
des catégories \unsure{grdes} devant $\pi$ (p. 81) : Que ces
\scanpage{263}\mspage{179}
% archive sheet number: 262
conditions impliquent $\deleted{M_\pi}\mathcal F^\pi=M_\pi$, i.e. $M_\pi\subset\mathcal F^\pi$,
dans le cas où $\pi>\pi_0$, se voit comme (4.7.9.1),
en utilisant le lemme 4.7.9.1.

\uline{Définition} 4.7.13 : $1^{\mathrm{o}}$) Un cardinal $\pi\geq\pi_0$ est
dit \inserted{\uline{cardinal}} \uline{utile} pour la fcga $(\mathcal F^\pi)_{\pi\geq\pi_0}$, \deleted{\unsure{si elle|s'il}}
s'il satisfait les conditions équivalentes
a) et b) de 4.7.12, et si de plus, dans
le cas où $\pi=\pi_0$, on a \deleted{(4.7.12.1)} $\mathcal F^\pi=M_\pi$. (\unsure{Donc|Dans}
si $\pi$ est un cardinal utile, on a \inserted{\unsure{tjrs}} \deleted{\unsure{4.7}} $\mathcal F^\pi=M_\pi$.)
\deleted{Donc si $\pi$ est utile, on a}
\deleted{On dit que $\pi$ est un card}
Dans le cas contraire, on dit que $\pi$ est
un \uline{cardinal superflu} (pour la fcga
donnée).

$2^{\mathrm{o}}$) Un cardinal $\pi$ est dit \uline{régulier}
pour la fcg donnée, \deleted{s'il satisfait} \inserted{s'il est utile}
\deleted{\unsure{à}} \inserted{et si de plus} la majoration de cardinaux énoncée dans
(4.7.7.1) est valable.

4.7.\unsure{14}.\deleted{\unsure{1}} Ainsi, tous les cardinaux $\pi\geq\pi_0$ tels que
$\pi^{\pi_0}=\pi$ sont utiles, et même réguliers pour
la filtration cardinale donnée. Notons
que si tous les cardinaux $\pi\geq\pi_0$ sont
utiles, alors on a $\mathcal F^\pi=M_\pi$ pour tout
$\pi\geq\pi_0$, donc la fcga est canonique. Notons
aussi que dans tous les cas, si $(\mathcal F^\pi)_{\pi\geq\pi_0}$ est
une fcga, alors $(M_\pi)_{\pi\geq\pi_0}$ en est une aussi.
\scanpage{264}\mspage{180}
% archive sheet number: 263
En effet, la condition a) des \unsure{fcg|feg} (cf 4.7.7) résulte de
la $=$ condition pour $(\mathcal F^{\unsure{\pi}})$, et de (4.7.10.1)
$M_\pi\subset\mathcal F^{\pi^{\pi_0}}$. La condition b) est une \deleted{\unsure{condition}} \inserted{\unsure{propriété}}
générale de la $\pi$-accessibilité, compte tenu
que $(\mathcal F^{\pi_0})_0$ est une petite sous-catégorie génératrice
formée d'objets $\pi_0$-accessibles de $M$. La condition \unsure{d|b}) est satisfaite tautologiquement, et
il reste à vérifier la condition c). Mais
on a vu que pour $\pi=\pi^{\pi_0}$, on a $M_\pi=\mathcal F^\pi$
(4.7.9.1), donc la condition c) pour $(M_\pi)$
résulte de la $=$ condition pour $(\mathcal F^\pi)$.

Ainsi, on trouve

\uline{Proposition} 4.7.15. Si $(\mathcal F^\pi)_{\pi\geq\unsure{\pi_0|\pi}}$ est une
fcga sur la $\unsure{U|\mathcal U}$-catégorie $M$, il en est
de $=$ de $(M_\pi)_{\pi\geq\pi_0}$. De plus \deleted{tout}
cardinal utile (resp. régulier) pour la
première filtration, l'est pour la \deleted{\unsure{f}} seconde
(appelée la \deleted{filtration cardinale} \inserted{fcga \uline{canonique}} \uline{associée}
à la filtration cardinale donnée).

La dernière assertion résulte du
fait que si $\pi$ est utile, pour $(\mathcal F^\pi)$, alors
on a $\mathcal F^\pi=M_\pi$..

4.7.16. Il me semble qu'on \unsure{peut} dire que la
fcg canonique associée à une fcga donnée
\scanpage{265}\mspage{181}
% archive sheet number: 264
est toujours ``meilleure'' que celle-ci (si celle-ci
\deleted{\unsure{lui}} n'est \deleted{\unsure{pas déjà}} déjà canonique, i.e. si les
deux ne sont égales), à cause des deux
faits suivants : a) \inserted{\unsure{surtout|d'abord}} parce qu'elle est \uline{canonique},
i.e. ne dépend d'aucun autre choix que
celui du seul cardinal $\pi_0$, et b) parce
qu'il y a a priori plus de cardinaux
utiles et de cardinaux réguliers pour
la fcg canonique. \inserted{Donc au total, il me semble que les filtrations cardinales \unsure{g.} les seules utiles sont celles qui sont \deleted{\unsure{access}} canoniques.}
Notons qu'une filtration cardinale (au sens
de SGA 4 \unsure{I \S|I}, que je conserve ici) \deleted{qui} est de plus
\uline{accessible} \inserted{(i.e. $\mathcal F^\pi\unsure{\subset|=}M_\pi$ pour tt $\pi\geq\pi_0$)}, \unsure{n'est|en est} autre qu'une \deleted{\unsure{filtr}}
fcga canonique, telle que \uline{tous} les
cardinaux sont réguliers (a fortiori utiles)
pour \unsure{ladite|la dite} filtration. Je présume \deleted{que}
(sans avoir cependant construit d'exemple)
que les filtrations cardinales canoniques
des catégories de la forme $M=\operatorname{Ind}_{\pi_0}(C)$,
$C$ une \inserted{petite} catégorie \uline{karoubienne} (pour
que $\pi_0$ soit un cardinal utile) \deleted{\unsure{don}}
ne sont pas forcément des fc (non généralisées, i.e. au sens SGA), i.e. je présume
qu'il arrive qu'il y ait des cardinaux $\pi\geq\pi_0$
non réguliers, et $=$ superflus.
\scanpage{266}\mspage{182}
% archive sheet number: 265
Il me faut enfin étudier ce qui se
passe quand on change le cardinal initial
$\pi_0$ d'une fcga :

\uline{Proposition} 4.7.17 Soit $(\mathcal F^\pi)_{\pi\geq\pi_0}$ une fcga
sur la $\unsure{U|\mathcal U}$-catégorie $M$, et soit $\pi'_0$ un
cardinal $\geq\pi_0$. Alors $(\mathcal F^\pi)_{\pi\geq\pi'_0}$ est une
fcga sur $M$, qui est canonique si
celle de départ \deleted{\unsure{l'}} est canonique. \unsure{Pour|Tout}
\inserted{\unsure{un}} cardinal \inserted{$\pi\geq\pi'_0$, \deleted{\unsure{qui est}} celui-ci est} utile, \deleted{(resp. régulier)} pour
$(\mathcal F^\pi)_{\pi\geq\pi'_0}$ ssi il l'est pour $(\mathcal F^\pi)_{\pi\geq\pi_0}$.
\inserted{Pour qu'il soit régulier pour $(\mathcal F^\pi)_{\pi\geq\pi'_0}$, il \uline{suffit} qu'il le soit pour $(\mathcal F^\pi)_{\pi\geq\pi_0}$, et c'est aussi nécessaire \unsure{(}}%
% margin note written sideways
\inserted{en supposant de plus que $2^\pi\geq\pi'_1\overset{\text{déf}}{=}\operatorname{Sup}(\pi'_0,\unsure{\text{Cardréd}|\mathrm{Cardrid}}\,\mathcal F^{\pi'_0})$, a fortiori \unsure{quand|que} $\pi\geq\pi'_1$.}

\uline{Dém} Les conditions a) b) sont trivialement
satisfaites pour la deuxième filtration,
\deleted{l'étant} \deleted{satisfaites} \inserted{\unsure{pour}} pour la première,
ainsi que d) si elle l'est pour la première.
Il en est de même de la condition
c) si $\pi=\pi^{\pi'_0}$, qui implique en effet
$\pi^{\pi_0}=\pi$, puisqu'en général
\[
\pi\leq\pi^{\pi_0}\leq\pi^{\pi'_0}\mathpunct{.}
\]

Reste à examiner la majoration des
cardinaux
\[
\begin{aligned}
\operatorname{card} I\leq\pi'^{\pi}\qquad &\text{si}\quad F\in\operatorname{Ob}\mathcal F^{\pi'},\\
&\text{\uuline{et}}\quad\pi'^{\pi}\geq\pi'_1,
\end{aligned}
\]
où cette fois $\pi'_1=\operatorname{Sup}(\pi'_0,\text{card réd}\,\mathcal F^{\pi'_0})$. \deleted{\unsure{Comme}}
\unsure{En} tout état de cause, on a $\pi'_1\geq\pi_1$, donc
\scanpage{267}\mspage{183}
% archive sheet number: 266
$\pi'^{\pi}\geq\pi'_1$ implique $\pi'^{\pi}\geq\pi_1$, donc on peut appliquer
la condition c) des \unsure{fcg|feg} pour $(\mathcal F^\pi)_{\pi\geq\pi_0}$ pour
conclure qu'elle reste encore vérifiée. Donc
en somme, \uline{toutes} les conditions des fcg
(resp. fcga) sont trivialement vérifiées pour
la fcg \uline{restreinte} de la première.

De $=$, le fait qu'un cardinal $\pi'\inserted{\geq\pi'_0}$ est
utile \deleted{pour} est trivialement indépendant de
quelle des deux filtrations on considère,
elle ne dépend que de la donnée du
cardinal $\pi'$ et de la sous-catégorie $\mathcal F^{\pi'}$
de $M$. Reste la majoration des cardinaux, dans le cas où on s'intéresse à
la \uline{régularité} de $\pi'$ pour les \deleted{\unsure{filtr}} fcga
envisagées. La majoration exigée est
la $=$, mais la condition \inserted{\deleted{\unsure{card}} préalable} de
validité de cette majoration, savoir
\[
\pi'^{\pi}\geq\pi'_1\mathpunct{.}\qquad\text{resp.}\quad\pi'^{\pi}\geq\pi_1
\]
n'est pas la même. Comme $\pi'_1\geq\pi_1$,
on voit que la condition \unsure{pré=alable} est
moins restrictive dans le cas de la
première filtration, que de la seconde, d'où
le ``il suffit'' dans la dernière assertion. Mais
\uline{lorsque} $\deleted{\unsure{\pi|2}}2^\pi\geq\pi'_1$, \deleted{\unsure{or}} la plus restrictive est automatique%
\scanpage{268}\mspage{184}%
% archive sheet number: 267
ment vérifiée, puisque $\pi'^\pi\geq2^\pi\geq\pi'_1$, et alors
la \deleted{validité} régularité de $\pi$ pour la \unsure{2ème}
filtration implique la régularité pour la
première, \unsure{cqfd|cqf}.

\uline{Corollaire} 4.7.18. \inserted{Supposons $\pi_0$ régulier pour la fcga$^{*}$.} Si on restreint la fcga $(\mathcal F^\pi)_{\pi\geq\pi_0}$
% slanted
\margin{$^{*}$ Il suffit pour cela que $\pi_0$ soit utile, i.e. $M\xleftarrow{\approx}\operatorname{Ind}_{\pi_0}(\mathcal F^{\pi_0}=M_{\pi_0})$, \deleted{\unsure{cf}} cela implique régulier, cf 4.8.5.}
aux cardinaux \inserted{$\pi\geq\pi'_0$, où} $\pi'_0\geq\pi_1$, \deleted{est régulier pour la fcg} alors on trouve
une fcga où la majoration des cardinaux
\deleted{\unsure{est}} (pour $\pi=\pi^{\pi_0}$) est valable \deleted{soit la condition
(4.7.7.1) est valable} sans condition préalable
$\pi'^\pi\geq\pi'_1$ ; plus précisément, on a alors
\[
\pi'_1\leq\deleted{\unsure{2^{\pi'}\pi_0}}2^{(\pi'_0{}^{\pi_0})}
\quad(\text{a fortiori }\pi'_1\leq2^{(2^{\pi'_0})}),
\tag*{(4.7.18.1)}
\]
donc la condition préalable $\pi'^\pi\geq\pi'_1$ est
automatiquement vérifiée \inserted{si $\pi\geq\pi'_0$, $\pi=\pi^{\pi'_0}$, $\pi'>\pi$}
(cf. remarque 4.7.7.2).

Ceci \deleted{résulte du résultat de majoration} \inserted{\unsure{s'aperçoit|s'ensuit} \unsure{avec|sur} la majoration}
de cardinaux suivante (appliquée au cas de $\pi=\pi'_0$)

\deleted{Corollaire}
\uline{Lemme} 4.7.19 : Supposons $\pi_0$ régulier pour la
fcga $(\mathcal F^\pi)_{\pi\geq\pi_0}$. Alors pour tt cardinal $\pi\geq\pi_1$
$(\overset{\text{déf}}{=}\operatorname{Sup}(\pi_0,\text{Card réd}\,\mathcal F^{\pi_0}))$, on a
\[
\mathrm{Card\,red}\,\mathcal F^\pi\leq2^{(\pi^{\pi_0})},
\tag*{(4.7.19.1)}
\]
et si $\pi'>\pi$, \deleted{et} $F\in\operatorname{Ob}\mathcal F^{\pi'}$, on a
\deleted{Démonstration de}
\[
\text{card réd}\,\mathcal F^{\unsure{\pi|\pi'}}_{/F}\leq\pi'^{(\pi^{\pi_0})}.
\tag*{(4.7.19.2)}
\]
% slanted
\margin{\uline{NB} J'ignore si on peut \unsure{remplacer} \deleted{\unsure{pour}} dans ces majorations $\pi^{\pi_0}$ par $\pi$.}
\scanpage{269}\mspage{185}
% archive sheet number: 268
\uline{Dém du lemme}. Comme $\pi_0$ est régulier,
posant $C=\mathcal F^{\pi_0}=M_{\pi_0}$, on trouve $M\approxeq\operatorname{Ind}_{\unsure{\pi_0}}(C)$,
donc \unsure{OPS} $M=\operatorname{Ind}_{\pi_0}(C)$, et appliquer les
résultats prouvés pour \unsure{les catégories|une catégorie} de cette
forme. Comme $\mathcal F^\pi\subset M^{\deleted{\unsure{\pi}}}_{\pi}$, \inserted{$\mathcal F^{\pi'}\subset M_{\pi'}$}, les majorations
(4.7.9.1) \inserted{et (4.7.9.2)} sont contenues ds
\[
\left.\left\{\begin{aligned}
\operatorname{card\,red}M_\pi &\leq 2^{(\pi^{\pi_0})}\\
\operatorname{card\,\text{réd}}M_{\pi/F} &\leq \pi'^{(\pi^{\pi_0})}
\quad\text{si }F'\in\operatorname{Ob}M_{\pi'},\\
&\hspace{4em}\text{avec }\pi'\geq\pi\geq\unsure{\pi_{0,1}|\pi_1}
\end{aligned}\right.\right\}\ \text{pour }\pi\geq\pi_1 .
\]
% slanted
\margin{Comparer aux majorations $\left\{\begin{aligned}\operatorname{Card\,red}\mathrm{Filt}^\pi&\leq2^\pi\\\operatorname{Card\,red}\mathrm{Filt}^{\pi}_{/F}&\leq\deleted{\unsure{\pi}}\pi'^{\pi}\end{aligned}\right.$ (4.7.\unsure{4|9}.1) et (4.7.\unsure{4|9}.2) si $F\in\operatorname{Ob}\mathrm{Filt}^{\pi'}$, avec $\pi'\geq\pi\geq\pi_1$, majorations plus \unsure{agréables|exigeantes} que celles de (4.7.\unsure{19|15}). Il faudrait quand $=$ s'assurer sur un exemple \deleted{que} si \unsure{les|ces} majorations plus \unsure{symp.|exigeantes} ne \unsure{st} pas valables dans le cas général.}
La première de ces majorations n'est autre que
la deuxième majoration dans
4.7.\unsure{4|7}.1. Notons qu'on a bien
\[
2^{(\pi^{\pi_0})}\leq2^{(2^\pi)},
\]
résulte en effet de
\[
\pi^{\pi_0}\leq2^{\unsure{\pi_0|\pi}},
\]
qu'on prouve en notant que $\pi<2^\pi$,
d'où $\pi^{\pi_0}\leq(2^\pi)^{\pi_0}=2^{\pi\mathbin{.}\pi_0}=2^\pi$ puisque $\pi\geq\pi_0$.
Il reste à établir la deuxième majoration
\[
\begin{aligned}
\operatorname{card\,red}M_{\pi/F}&\leq\pi'^{(\pi^{\pi_0})}\quad\text{si }F\in\operatorname{Ob}M_{\pi'},\\
&\hspace{6em}\text{avec }\pi'\geq\pi\geq\pi_1 .
\end{aligned}
\tag*{(\unsure{$\alpha$|as|bis})}
\]
(On n'\deleted{\unsure{en}} a pas besoin pour 4.7.18, \deleted{\unsure{et|si}} mais
ce sera sans doute utile par ailleurs). On
utilise $M_{\pi'}\subset\mathrm{Filt}^{(\pi'^{\pi_0})}$ si $\pi'\geq\pi_1$
(4.7.6.2), donc si $F\in\operatorname{Ob}M_{\pi'}$ on a
\[
\deleted{F=\varinjlim_{I}x_i}\ \Big|\ 
\deleted{\begin{aligned}&\unsure{\mathrm{les}}\ x_i\in C,\\
&I\ \text{grd devant }\pi_0,\\ &\operatorname{Card}I\leq\pi^{\pi_0},\end{aligned}}
\]
\scanpage{270}\mspage{186}
% archive sheet number: 269
\[
\operatorname{card}F(x)\leq\pi'^{\pi_0}\qquad\forall\ x\in\operatorname{Ob}C,
\]
donc si $G\in\operatorname{Ob}M_\pi\subset\operatorname{Ob}\mathrm{Filt}^{\pi^{\pi_0}}$, \deleted{alors}
\[
\operatorname{card}G(x)\leq\pi^{\pi_0}\qquad\forall x\in\operatorname{Ob}C,
\]
on trouve
\[
\operatorname{card}\operatorname{Hom}(G,F)\leq
\underbrace{\underset{\leq\pi_1}{\tikz[baseline=(q.base)]{\node[inner sep=0pt,outer sep=0pt] (q) {$\operatorname{card}(\operatorname{Ob}C)$};\draw ([yshift=-.3pt]q.south west) -- ++(0,-2.2pt) -| ([yshift=-.3pt]q.south east);}}
\mathbin{.}(\pi'^{\pi_0})^{(\pi^{\pi_0})}}_{
\pi'^{\pi_0\mathbin{.}\pi^{\pi_0}}=\pi'^{(\pi^{\pi_0})}\mathpunct{;}}\mathpunct{,}
\]
\deleted{d'où}

d'où
\[
\operatorname{card}\operatorname{Ob}M_{\pi/F}\leq
\underbrace{\underset{\leq2^{(\pi^{\pi_0})}}{\tikz[baseline=(q.base)]{\node[inner sep=0pt,outer sep=0pt] (q) {$\operatorname{card}(\operatorname{Ob}M_\pi)$};\draw ([yshift=-.3pt]q.south west) -- ++(0,-2.2pt) -| ([yshift=-.3pt]q.south east);}}
\mathbin{.}\pi'^{(\pi^{\pi_0})}}_{
\deleted{\unsure{\gamma}}\gamma=\pi'^{(\pi^{\pi_0})}},
\]
et enfin
\[
\operatorname{card}\mathrm{Fl}\,M_{\pi/F}\leq
\underbrace{\underset{\gamma=\pi'^{(\pi^{\pi_0})}}{
\tikz[baseline=(q.base)]{\node[inner sep=0pt,outer sep=0pt] (q) {$\deleted{\unsure{\gamma}}\gamma\mathbin{.}\gamma$};\draw ([yshift=-.3pt]q.south west) -- ++(0,-2.2pt) -| ([yshift=-.3pt]q.south east);}}\mathbin{.}
\underset{\leq2^{(\pi^{\pi_0})}}{\tikz[baseline=(q.base)]{\node[inner sep=0pt,outer sep=0pt] (q) {$\operatorname{card}(\mathrm{Fl}\,M_\pi)$};\draw ([yshift=-.3pt]q.south west) -- ++(0,-2.2pt) -| ([yshift=-.3pt]q.south east);}}}_{
\pi'^{(\pi^{\pi_0})}},
\]
ce qui est la majoration annoncée
(la $=$ que celle prouvée dans 4.5.\unsure{17|12}, la
dém. étant ess. la même).
\scanpage{271}\mspage{\unsure{187}}
% archive sheet number: 270
4.7.20. Une \unsure{fcg|feg} canonique est dite \uline{maximale}
\inserted{(ou \uline{strictement canonique})},
si elle n'est pas \deleted{\unsure{déduit}} déduite d'une fcg canonique
str. plus grande $(\mathcal F^\pi)_{\pi\geq\pi_{00}}$ (donc avec
$\pi_{00}<\pi_0$). \deleted{plus} Donc si $M$ admet une
fcga, alors elle admet même une
fcg str. canonique, et cette dernière est de
plus unique. \deleted{\illegible{3}} On va revenir
là-dessus dans la section suivante.

4.8. \uline{Catégories} $\unsure{U|\mathcal U}$-\uline{accessibles}.

\deleted{\illegible{1}} \inserted{\uline{Proposition}} 4.8.1. Soit $M$ une $U$-catégorie,
\inserted{$\pi_0$ un cardinal infini. Considérons}
\deleted{\unsure{a)}} les conditions suivantes : \deleted{sont équivalentes :}

\deleted{\unsure{b)}}a) $M$ satisfait \unsure{à|:} $L\pi_0$ (i.e. est stable par \inserted{petites} $\varinjlim$
\unsure{grdes} devant $\pi_0$), et contient une petite
sous-catégorie pleine $\unsure{S|\mathcal S}$ \uline{strictement
génératrice}, formée d'objets $\pi_0$-accessibles.

b\deleted{\unsure{c}}) Il existe une petite catégorie $C$,
et une \deleted{\unsure{foncteur} pl. fidèle \unsure{non accessible|non essentiel}} \inserted{équivalence de catégories}
\[
M\approx\operatorname{Ind}_{\pi_0}(C)
\]

\deleted{\unsure{b)}}c) $M$ admet une filtration cardinale
généralisée accessible $(\mathcal F^\pi)_{\pi\geq\pi_0}$, ayant
$\pi_0$ comme cardinal initial, et telle que
$\pi_0$ soit un cardinal utile pour ladite
filtration.
\scanpage{272}\mspage{188}
% archive sheet number: 271
c') Comme c), \deleted{\unsure{mais}} avec $\pi_0$ cardinal
régulier (et non seulement utile) pour la
filtration card. gén. \inserted{acc.} en question, et cette
\deleted{\unsure{la}} filtration étant supposée de plus \uline{canonique}.
% margin note written sideways
\margin{On a alors le diagramme d'implications
\diagram{\begin{tikzcd}[ampersand replacement=\&,column sep=3pt,cells={nodes={inner sep=1pt,font=\scriptsize}}]
c'\arrow[r,Leftrightarrow]\&c\arrow[r,Leftrightarrow]\&b\arrow[r,Rightarrow]\&a
\end{tikzcd}}}

4.8.1.1. Il me faut ici préciser la
terminologie, sur ce que j'entends par
sous-catégorie (pleine) $C$ \uline{strictement} \uline{génératrice}
dans $M$. Pour ceci, je me réfère à
SGA 4 I 7.2, (i), aux conditions équivalentes b) et c) \unsure{du|de} loc. cit.

\textcircled{b} $\forall F\in\operatorname{Ob}M$, la flèche canonique
\[
\varinjlim_{i\in C/F}x_i\longrightarrow F
\]
est un iso ;

\textcircled{c} Le foncteur canonique
\[
\begin{aligned}
M&\longrightarrow C^\wedge\\
F&\longmapsto\bigl(\underset{\substack{\unsure{\rotatebox{90}{\(\in\)}}\\\operatorname{Ob}C}}{x}
\unsure{\mapstochar\relbar|\mapsto}\underset{\in\operatorname{Ob}(\mathrm{Ens})}{\underline{\operatorname{Hom}_M(x,F)}}\bigr).
\end{aligned}
\]
est pleinement fidèle.

Ces conditions impliquent

a) $C$ est génératrice \inserted{dans $M$} par épimorphismes stricts, \inserted{i.e. pour tt $F$ \unsure{dans|de} $M$, la famille des flèches $x\to F$ de source dans $C$, de but $F$, est épimorphique stricte.}
et lorsque $M$ est stable par produits
fibrés, cette condition a) est équivalente
aux précédentes. Dans loc. cit. j'affirme
\scanpage{273}\mspage{189}
% archive sheet number: 272
cette équivalence sans hypothèse restrictive sur
$M$ --- mais en annotation marginale
j'affirme que cette assertion est \uline{fausse},
en dehors de l'hypothèse d'existence des
produits fibrés dans $M$. Pourtant, je
viens de vérifier la démonstration donnée
dans loc. cit de l'implication $a\Longrightarrow b$,
et elle est correcte. Mais peut-être l'annotation résulte-t-elle d'une confusion
entre la condition a), et \deleted{\unsure{la}} la condition suivante, plus faible a priori et qui
lui ressemble comme un frère :

a') Il existe une famille de flèches
\[
u_i:\unsure{x_i|X_i}\longrightarrow F,
\]
de source dans $C$ et
de but $F$, qui est \unsure{épimorphique|épimorphiques}
stricte (cf SGA 4 I 10.3).

On a envie de dire que quand on \uline{agrandit}
une famille de flèches \unsure{épimorphique|épimorphiques}
stricte, en \unsure{rajoutant|ajoutant} \deleted{\unsure{des}} d'autres flèches de
$=$ but, elle doit rester épim.
stricte (``a fortiori'') --- ce qui donne
l'impression que a') doit impliquer \uline{a)},
d'où leur \deleted{\unsure{str.}} équivalence. Mais si on \unsure{essaye|essaie} de prouver cette propriété hypothétique \deleted{pour}
\scanpage{274}\mspage{190}
% archive sheet number: 273
de stabilité pour les familles épimorphiques
strictes, on voit qu'on n'y arrive que
si on suppose que la famille initiale
est de plus \uline{épimorphique} \uline{universelle}
(si les produits fibrés existent dans $M$),
ou (sans la \unsure{réserve|dernière} \unsure{précédente|précision}) si \deleted{\unsure{les}} pour
tt $\deleted{\unsure{x}}F'\longrightarrow F$, le crible de $\deleted{\unsure{C}}M_{/F'}$ image
inverse du crible de $M_{/F}$ engendré par
les $x_i\unsure{\mathord{-}|\to}F$ est un ``crible épimorphique''.
\deleted{\unsure{Cela}} Cela fait entrevoir en effet qu'il doit
y avoir des cas où la condition a') est
satisfaite, mais non la condition a)
(équivalente à b) et à c)).

Mais nous sommes alors confrontés
à la situation très déplaisante : si $C$ est
une sous-catégorie strictement génératrice de
$M$, il ne s'ensuit pas pour autant qu'une
sous-catégorie pleine $C'$ \deleted{\unsure{par}} \inserted{de $M$} telle que $C'\supset C$
soit également strictement génératrice !
Or j'ai toujours utilisé \inserted{\unsure{précédemment|précisément}}, sans y réfléchir,
deux fois, que si $C$ est \deleted{\unsure{g}} str. génératrice,
il en est de $=$ des $C'\supset C$. Donc il
faudra que je sois à présent très attentif,
\scanpage{275}\mspage{191}
% archive sheet number: 274
et que je vérifie si cette erreur \unsure{entache|entâche}
\deleted{\unsure{entache d'erreur}} rend inopérantes certaines
démonstrations, voire même, rend faux certains
énoncés. \unsure{Aussi|Ainsi}, il me faudrait donner
un exemple où $C$ est str. gén., mais
$C'\supset C$ ne l'est pas.

Finalement, réflexion faite, ces ennuis \unsure{craints|criants}
disparaissent. Je vais énoncer un résultat
en forme :

\deleted{Lemme 4.8.1}
\uline{Proposition} 4.8.2. Soient $C,M$ deux \inserted{$\unsure{U|\mathcal U}$-}catégories,
$i:C\unsure{\longrightarrow|\Longrightarrow} M$ un foncteur, d'où un
foncteur
\[
\left|\begin{aligned}
\varphi &: M\longrightarrow C^\wedge\\
\varphi(\unsure{X|x})&=(x\longmapsto\operatorname{Hom}_M(i(x),X)
\end{aligned}\right.
\tag*{(4.8.2.1)}
\]
(foncteur composé $M\xrightarrow{\mathrm{can}}M^\wedge
\xrightarrow[\text{i.e. composition avec }i]{i^*}C^\wedge$).

\deleted{$1^{\mathrm{o}}$) Considérons les conditions suivantes}
Pour tout $X$ dans $M$, on considère la
\deleted{\unsure{a) l}} catégorie $C_{/X}\simeq C_{/\varphi(\unsure{x|X})}$, et la flèche
canonique dans \deleted{\unsure{la}} $M^\vee$
\[
\varinjlim_{x\in C_{/X}} i(x)\longrightarrow X\mathpunct{.}
\tag*{(4.8.2.2)}
\]

$\unsure{1|2}^{\mathrm{o}}$) Considérons les conditions suivantes

\textcircled{a} \inserted{Pour tt $X\in\operatorname{Ob}M$,} la famille des flèches $i(x)\longrightarrow X$, indexée
par $\operatorname{Ob}C_{/X}$, est épimorphique stricte.
\scanpage{276}
\mspage{192}
% archive sheet number: 275

\textcircled{b} Pour tt $X\in\mathrm{Ob}\,M$, la flèche can. 4.8.2.2.\inserted{est iso,}
\inserted{i.e.} fait de $X$ une $\varinjlim$ du foncteur
canonique $C_{/\unsure{X|x}}\longrightarrow M$.

\textcircled{c} Le foncteur $\varphi$ (4.8.2.1) est pleinement
fidèle.

On a alors les implications
\diagram{\begin{tikzcd}[ampersand replacement=\&]
 b \arrow[r,Leftrightarrow] \& c \arrow[r,Rightarrow] \& a,
\end{tikzcd}}
et \deleted{\unsure{si|s}}si on suppose $i$ pl. fidèle, alors les
trois conditions $a,b,c$ sont équivalentes.

[\uline{NB} On dit que le foncteur $i:C\longrightarrow M$
est \uline{strictement générateur}, si les conditions
équivalentes $b,c$ sont satisfaites. Si $C$ est
une sous-catégorie de $M$ (\unsure{pas|non} néc. pleine),
on dit que $C$ est \uline{strictement génératrice}
si le foncteur d'inclusion est strictement
générateur, i.e. satisfait aux conditions
équivalentes $b)$, $c)$. Pour \deleted{$C$ est une \unsure{sous-c|sous-c.}}
\inserted{\unsure{ce,|en,|oui,} il faut donc}
que la condition $a)$ soit satisfaite i.e.
que pour tt $X$ dans $M$ la famille
des flèches $x\longrightarrow X$ de source dans
$C$, de but $X$, soit épimorphique stricte,
et cette condition est aussi suffisante
quand $C$ est sous-catégorie pleine de $M$.]

\margin{descend après $2^{\mathrm o}$}
$3^{\mathrm o})$ Considérons un diagramme commutatif
de foncteurs de $U$-catégories\,:
